\section{Offset geometry and residual relation spaces}
\label{p1:offset:section}

The group budget used in the main proof keeps the product of primes
sharing an offset. Two further geometric consequences explain what such
products can force. A line of offsets has a bounded total prime product;
a large product on three offsets gives a plane involving all three
windows. These statements also give a more detailed description of the
relation space after a covering set of coordinates has been removed.

Throughout this section $B=L+1$, $I_N=[N,2N)\cap\mathbb Z$, and
$H=C_0B$ for a fixed $C_0>1$. We consider three distinct regular windows
with starts $X=(x_1,x_2,x_3)\in I_N^3$ and pairwise direct relative
spaces. Their offsets lie in
$\mathcal O=\{-1,0,\ldots,L-1\}$, whose diameter is less than $B$.
Every prime $\ell>H$ having odd valuation in all three windows has a
unique offset $u_\ell\in\mathcal O^3$. Such primes are called shared
in the statements below. No independence of the three-window system
is assumed. Whenever a primitive mass is mentioned, it is the sum of
the original affine joint probabilities over ordered primitive triples.
Each of these probabilities is at most $p^2$.

\subsection{Whole prime products on lines and planes}

\begin{lemma}[The prime product on an offset line]
\label{p1:offset:line-product}
Let $\mathcal L$ be a line in offset space, and let
$\mathcal Q_{\mathcal L}$ be the shared primes whose offsets lie on it.
If their distinct offsets consist of one point, then
\[
 \prod_{\ell\in\mathcal Q_{\mathcal L}}\ell\le3N.
\]
If there are at least two distinct offsets, the line has a primitive
integer direction $d$ and
\begin{equation}
 \prod_{\ell\in\mathcal Q_{\mathcal L}}\ell
       \le6N\|d\|_\infty\le6NB.
 \label{p1:offset:line-product-bound}
\end{equation}
For $Y\ge H$, the number of distinct offsets on this line supplied by
primes greater than $Y$ is at most
\[
 M_Y=1+\frac{\log(6NB)}{\log Y}=O_{C_0,C}(B/\log B)=o(B).
\]
\end{lemma}
\begin{proof}
At a single offset the product of the distinct primes divides each
coordinate of $X+u$, all of which are at most $3N$ for large $N$.
Otherwise fix one occurring offset $u_0$. Every integer point on the
line has the form $u_0+td$, with $t\in\mathbb Z$. Each of its shared
primes divides the same integer vector
\[
 d\times(X+u_\ell)=d\times(X+u_0).
\]
This vector is nonzero. If it were zero, positivity of $X+u_0$ would
make all coordinates of $d$ nonzero. Two distinct offsets would give
two positive proportional value vectors, with two distinct values in
each block. Their projections on any two blocks would give equal
nonzero square classes of even selections. Injectivity and the direct
sum of those two spaces exclude this relation.

A nonzero coordinate of the displayed cross product has absolute value
at most $6N\|d\|_\infty$ and is divisible by every prime in the
product. Two offsets in $\mathcal O^3$ give $\|d\|_\infty<B$.
This proves \eqref{p1:offset:line-product-bound}. Choose one prime per
distinct offset and take logarithms to obtain the last assertion;
the extra one covers a line carrying only one offset.
\end{proof}

\begin{lemma}[An integer product tail]
\label{p1:offset:integer-tail}
For fixed $s>1$ and $T\ge1$,
\[
 \sum_{\substack{a,b\ge1\\ab>T}}(ab)^{-s}
       \ll_s T^{1-s}\log(2T).
\]
\end{lemma}
\begin{proof}
For $a\le T$, the inner sum is at most
$C_s(T/a)^{1-s}$, giving
$C_sT^{1-s}\sum_{a\le T}a^{-1}$ after multiplication by $a^{-s}$.
For $a>T$, use the full sum $\zeta(s)$ and
$\sum_{a>T}a^{-s}=O_s(T^{1-s})$.
\end{proof}

Let $\mathscr P_B$ be the family of distinct affine planes
\[
 A\cdot X=h,\qquad
 A=(u_2-u_1)\times(u_3-u_1)\ne0,
 \qquad h=-A\cdot u_1,
 \qquad u_1,u_2,u_3\in\mathcal O^3.
\]
There are at most $B^9$ such planes, and these representations satisfy
$\|A\|_\infty\le2B^2$ and $|h|\le6B^3$.

\begin{proposition}[A forced plane uses all three coordinates]
\label{p1:offset:forced-plane}
Two distinct shared primes satisfy $\ell\ell'\le6NB$.
If three distinct shared primes have product greater than $18NB^2$,
then $X$ belongs to a plane of $\mathscr P_B$ whose normal satisfies
$A_1A_2A_3\ne0$. Every further shared prime $\ell>12B^3$ has
\begin{equation}
 A\cdot u_\ell=-h.
 \label{p1:offset:subsequent-offset}
\end{equation}
The primitive mass of triples belonging to two distinct planes of
$\mathscr P_B$ is $O_C(N^{-1+o(1)})$.
\end{proposition}
\begin{proof}
For two primes, either they share a value in one of two selected blocks,
so their product is at most $3N$, or their occurrences are distinct in
both. In the latter case their product divides
$nm'-n'm$. Its vanishing gives the forbidden pair relation
$nn'mm'=(nm')^2$. Since $n,m\le3N$ and both within-block
differences have absolute value less than $B$, its nonzero absolute
value is at most $6NB$.

For three primes, their product divides
\[
 \det(X+u_1,X+u_2,X+u_3)
 =(X+u_1)\cdot((u_2-u_1)\times(u_3-u_1)),
\]
whose absolute value is at most $18NB^2$. Hence it vanishes.
The offsets cannot lie on a line, by
Lemma~\ref{p1:offset:line-product}, including the case of repeated
offsets. Thus $A\ne0$ and $A\cdot X=-A\cdot u_1$.

Suppose, for example, that $A_1=0$. The other two coefficients must
both be nonzero, since $A\cdot(X+u_i)=0$ cannot have just one
nonzero coefficient at a positive vector. The projections of all three
value vectors onto coordinates two and three are proportional. If two
projections differ, both their coordinates differ, giving a forbidden
pair relation. If none differ, the two offset differences are parallel
to the first axis, contradicting $A\ne0$. The other coordinates are
handled in the same way.

A further prime divides $A\cdot(X+u_\ell)=h+A\cdot u_\ell$,
an integer of absolute value at most $12B^3$. A larger prime forces
it to vanish. Finally, two distinct planes are disjoint or meet in a
line, which has at most $N$ integer points in $I_N^3$. There are at
most $B^{18}$ pairs of planes. Multiplication by the upper probability
$p^2$ gives the mass assertion.
\end{proof}

\begin{lemma}[Many planar supports force a small normal]
\label{p1:offset:normal-budget}
Let $A\cdot u=-h$ be a plane of offsets, with $A\in\mathbb Z^3$
primitive and nonzero, and $h\in\mathbb Z$. Let $m_{A,Y}$ count its
distinct offsets supplied by shared primes greater than $Y\ge H$.
Then
\begin{equation}
 m_{A,Y}\le M_Y+\frac{6B^2}{\|A\|_\infty}.
 \label{p1:offset:normal-count}
\end{equation}
If the hyperedges given by these offsets require at least $\delta L$
vertices to cover them, where $\delta>0$ is fixed, then for large $N$,
\[
 \|A\|_\infty\le24B/\delta,
 \qquad |h|\le72B^2/\delta.
\]
There are consequently $O_\delta(B^5)$ possible such integer planes.
\end{lemma}
\begin{proof}
The lattice $\Lambda=A^\perp\cap\mathbb Z^3$ has covolume
$\|A\|_2$. Indeed, a lattice basis $d,e$ can be completed by an
integer vector $w$ with $A\cdot w=1$ to a basis of $\mathbb Z^3$,
so $d\times e=\pm A$.
Choose a shortest nonzero lattice vector $d$, of length $\ell$.
It is primitive and extends to a lattice basis. Unless the cube has
at most one integer point on the plane, $\ell\le\sqrt3B$.
Parallel lattice lines in direction $d$ have spacing $\|A\|_2/\ell$.
At most $1+\sqrt3B\ell/\|A\|_2$ meet the cube, and each contains
at most $1+\sqrt3B/\ell$ integer points. By
Lemma~\ref{p1:offset:line-product}, each carries at most $M_Y$ supports.
It follows that
\[
 m_{A,Y}\le
 \left(1+\frac{\sqrt3B\ell}{\|A\|_2}\right)
       \min\left(M_Y,1+\frac{\sqrt3B}{\ell}\right)
 \le M_Y+\frac{6B^2}{\|A\|_\infty}.
\]
The omitted one-point case satisfies the same bound.
The covering number is no larger than the number of supports.
Since $M_Y=o(B)$ and $L=B-1$, a covering number at least $\delta L$
implies the asserted bound for $A$, with room in the constant.
The existence of an offset gives $|h|\le3B\|A\|_\infty$.
Counting the normals and the constants gives $O_\delta(B^5)$ planes.
\end{proof}

\subsection{Root separation and counting on a plane}

\begin{lemma}[At most one common root]
\label{p1:offset:root-separation}
Let $y=\alpha x+\beta$, with rational $\alpha\ne0$, pass through
a positive pair of regular windows with direct relative spaces.
Choose offsets distinct within each block. The two root lists
$-i$ and $-(\beta+j)/\alpha$ have at most one common element.
If these offsets are partitioned among equations with disjoint supports,
discarding at most two equations leaves squarefree polynomials with
pairwise disjoint root sets after substitution of the line.
\end{lemma}
\begin{proof}
A common root says $\beta+j=\alpha i$, and hence
$y+j=\alpha(x+i)$. Positivity makes $\alpha>0$.
Two common roots would give two distinct values on each side in the
same ratio. Their even square classes are nonzero by injectivity and
equal, contradicting directness. The sole possible common root belongs
to at most two equations. Removing those equations removes both
repeated roots within an equation and roots shared between equations.
\end{proof}

\begin{proposition}[Many transverse square products on prescribed planes]
\label{p1:offset:planar-components}
Fix $\eta,\varepsilon>0$ and $K>0$. Let a family of at most $B^K$
planes have integer equations
$A_1x+A_2y+A_3z=h$, with nonzero normal and all coefficients of
absolute value at most $B^K$.
The number of primitive triples on these planes admitting at least
$\eta B$ disjoint sets of values, each meeting all three windows and
having square product modulo primes at most $H$, is
$O_{C,C_0,\eta,K,\varepsilon}(N^{1+\varepsilon})$.
Their primitive mass is $O(N^{-1+\varepsilon})$.
\end{proposition}
\begin{proof}
A plane with only one nonzero coordinate coefficient cannot meet the
dyadic region for large $N$. Otherwise permute coordinates so that
$A_1A_2\ne0$. The budget of $3B$ values leaves at least $\eta B/2$
sets of size at most $D_0=\lceil6/\eta\rceil$.
Select $k+2$ of them, where $k$ is fixed below. Their equations are
\[
 w_\nu^2=d_\nu P_\nu(x)Q_\nu(y)R_\nu(z),
\]
with nonempty split products, disjoint root lists in each window,
total degree at most $D_0$, and squarefree $H$-smooth $d_\nu$.
Their offsets and coefficients have $N^{o(1)}$ possibilities for fixed
$k,D_0$, and $d_\nu\le\prod_{\ell\le H}\ell=N^{O_{C,C_0}(1)}$.

Fix $z$. The plane gives
$y=\alpha x+\beta$, where
$\alpha=-A_1/A_2$ and $\beta=(h-A_3z)/A_2$.
Lemma~\ref{p1:offset:root-separation} leaves $k$ equations whose
polynomials $P_\nu(x)Q_\nu(\alpha x+\beta)$ have simple disjoint
roots. Their square classes are geometrically independent, since
valuation at a root belonging to one equation detects that equation.
They define a geometrically integral curve of degree at least $2^k$
in the coordinates $(x,w_1,\ldots,w_k)$: its generic projection to
the $x$-line has that degree. Its degree is also bounded above in terms
of $k,D_0$ only.

Clearing powers of $A_2$ from the equations preserves the original
integer coordinates $w_\nu$. Their heights are at most $N^{q_0}$
for fixed $q_0=q_0(C,C_0,D_0,K)$ independent of $k,z$ and the motifs;
this follows from the original three-variable products.
Pila's uniform curve estimate \cite[Theorem~A]{pila1995} therefore
gives $O(N^{q_0/2^k+q_0\zeta})$ points for each motif and fixed $z$.
Choose $k$ sufficiently large, then $\zeta$ sufficiently small, to
make this exponent less than $\varepsilon/3$. Sum the $O(N)$ values
of $z$, the $N^{o(1)}$ motifs and the $B^K$ planes. Finally use
$q(\{x,y,z\})\le p^2$ to obtain the mass bound.
\end{proof}

\subsection{Deleting a cover of the three-block constraints}

For a separated primitive triple, let $\mathcal H_3$ have the $3B$
formal window positions as vertices and the three occurrences of each
shared prime as a hyperedge. Repeated supports may be identified.
Let $D$ be any vertex cover of these edges, with $t=|D|$.
For the relative relation space $R$, define
$R_D=\{v\in R:v|_D=0\}$.

\begin{lemma}[Covering cost and free components]
\label{p1:offset:cover-rank}
One has
\[
 \dim R\le t+\dim R_D.
\]
After coordinates in $D$ are set to zero, the equations from primes
greater than $H$ form a graph with possibly pinned vertices. If $c_3(D)$
counts its unpinned components meeting all three blocks, then
\[
 \dim R_D\le c_3(D)+O_{C,C_0}(B/\log B),
\]
uniformly in $D$. The product of the values in every free component
is a square modulo primes at most $H$.
\end{lemma}
\begin{proof}
The first statement is rank--nullity for the coordinate restriction
$R\to\mathbb F_2^D$. Every row previously of weight three meets
$D$, so its remaining weight is at most two. A row of weight two
equates its coordinates, and a row of weight one pins its coordinate
to zero; pinning propagates through the component.

A free isolated vertex was already $H$-defective. Otherwise a large
prime incident to it would remain either an edge or a pin. There are
$O(B/\log B)$ such vertices by Lemma~\ref{p1:arith:defects}.
A nonsingleton component cannot be contained in one block because
edges join distinct blocks. For a fixed pair of blocks, there are at
most $\pi(H)+2$ free components meeting exactly that pair: otherwise
their small-prime valuation vectors, augmented by the two block
parities, have a nonzero zero-sum combination. Disjointness of the
components would then give a forbidden relative pair relation.
These components also number $O(B/\log B)$.
All remaining free components meet the three blocks and each has
one binary coordinate before the additional equations are imposed.
This proves the dimension estimate. Every large-prime row meets a
free component zero or two times, proving its product assertion.
\end{proof}

\begin{proposition}[A six-vertex budget for the residual space]
\label{p1:offset:six-vertex}
For every fixed $\varepsilon>0$, there is a set of primitive triples
of total mass $O_{C,C_0,\varepsilon}(N^{-1+\varepsilon})$ outside
which every cover $D$ simultaneously satisfies
\begin{equation}
 \dim R_D\le\frac{3B-|D|}{6}
           +O_{C,C_0}(B/\log B)+O_{C,C_0,\varepsilon}(1).
 \label{p1:offset:six-vertex-bound}
\end{equation}
For fixed $\eta>0$, triples admitting some cover with
$\dim R_D\ge\eta L$ number $O_{C,C_0,\eta,\varepsilon}(N^{3/2+\varepsilon})$.
On a prescribed family of planes as in
Proposition~\ref{p1:offset:planar-components}, their number is
$O(N^{1+\varepsilon})$ and their primitive mass is
$O(N^{-1+\varepsilon})$.
\end{proposition}
\begin{proof}
For any fixed positive type $(a,b,c)$ with a singleton, sufficiently
many disjoint square products of that type have $O(N^{1+\varepsilon})$
primitive hosts. Fix the third centre and apply the degree-one part
of Lemma~\ref{p1:arith:geometric} to the other two, putting the
singleton in its linear variable. The specialized coefficients have
fixed polynomial height. The original smooth coefficients and offsets
have $N^{o(1)}$ choices; summing the third centre proves the count.
Multiplication by $p^2$ gives mass $O(N^{-1+\varepsilon})$.

Apply this exclusion to the finitely many permutations of
$(1,1,1)$, $(2,1,1)$, $(3,1,1)$ and $(2,2,1)$.
Outside their union, the number $s$ of free three-block components
of size at most five is bounded, for every cover $D$.
All other three-block components use at least six surviving vertices,
and the small ones use at least three. Therefore
\[
 6c_3(D)-3s\le3B-|D|.
\]
Lemma~\ref{p1:offset:cover-rank} proves
\eqref{p1:offset:six-vertex-bound}.

If $\dim R_D\ge\eta L$, the same lemma supplies a positive
proportion of disjoint three-block square products. Select a fixed
number of bounded-size ones. Fixing one centre and using the general
two-parameter estimate of Lemma~\ref{p1:arith:geometric} gives
$N^{3/2+\varepsilon}$ hosts after the same coefficient and offset
enumeration. On the prescribed planes use instead
Proposition~\ref{p1:offset:planar-components}.
The covers are never enumerated: their existence supplies the disjoint
sets, of which only a fixed number are selected for counting.
\end{proof}
